% =====================================================================
%  APÉNDICES
% =====================================================================
\clearpage
\renewcommand{\seccionactual}{APÉNDICE A · PHP EN UNA PÁGINA}
\section*{Apéndice A · PHP en una página}
\addcontentsline{toc}{section}{Apéndice A · PHP en una página}

\begin{multicols}{2}
\small
\parte{Abrir, hablar, cerrar}
\begin{php}{básico}
<?php
echo 'Hola';          // escribe
$x = 5;               // variable
?>
<p><?= $x ?></p>      <!-- atajo -->
\end{php}

\parte{Textos}
\begin{php}{comillas y punto}
$n = 'Ana';
echo "Hola $n";       // Hola Ana
echo 'Hola $n';       // Hola $n
echo 'Hola ' . $n;    // Hola Ana
\end{php}

\parte{Decidir}
\begin{php}{if}
if ($edad >= 18) {
    echo 'Pasá';
} elseif ($edad >= 13) {
    echo 'Con un adulto';
} else {
    echo 'No';
}
\end{php}

\parte{Listas y fichas}
\begin{php}{arrays}
$frutas = ['manzana', 'pera'];
echo $frutas[0];      // manzana

$plato = ['nombre' => 'Saice',
          'precio' => 38];
echo $plato['precio']; // 38
echo count($frutas);   // 2
\end{php}

\parte{Repetir}
\begin{php}{foreach en HTML}
<?php foreach ($menu as $p): ?>
  <li><?= $p['plato'] ?></li>
<?php endforeach; ?>
\end{php}

\parte{Escuchar}
\begin{php}{formularios}
$a = $_GET['nombre'] ?? '';
$b = $_POST['clave'] ?? '';
if ($_SERVER['REQUEST_METHOD']
      === 'POST') { /* ... */ }
\end{php}

\parte{Escudo (siempre al mostrar)}
\begin{php}{htmlspecialchars}
function e($t) {
  return htmlspecialchars($t,
    ENT_QUOTES, 'UTF-8');
}
echo e($_GET['nombre'] ?? '');
\end{php}

\parte{Base de datos}
\begin{php}{PDO}
$pdo = new PDO(
 'mysql:host=db;dbname=fogon;charset=utf8mb4',
 'fogon', 'fogon123');

$filas = $pdo->query(
  'SELECT * FROM menu')->fetchAll();

$q = $pdo->prepare(
  'INSERT INTO visitas (nombre, mensaje)
   VALUES (?, ?)');
$q->execute([$nombre, $mensaje]);
\end{php}

\end{multicols}

% =====================================================================
\clearpage
\renewcommand{\seccionactual}{APÉNDICE B · GLOSARIO}
\section*{Apéndice B · Glosario}
\addcontentsline{toc}{section}{Apéndice B · Glosario}

\renewcommand{\arraystretch}{1.35}
\begin{longtable}{>{\leavevmode\sffamily\bfseries\color{fuegooscuro}}p{3.6cm} p{12.2cm}}
Apache & El servidor web: recibe las peticiones del navegador y entrega las páginas (el \emph{mozo}). \\
Array & Una lista de valores (\texttt{['a', 'b']}) o una ficha con nombres (\texttt{['plato' => 'Saice']}). \\
Cliente & Quien pide la página: el navegador. \\
Concatenar & Pegar textos con el punto: \texttt{'Hola ' . \$nombre}. \\
Consulta preparada & Una consulta SQL con huecos (\texttt{?}) donde los datos del usuario van por separado. Protege contra la inyección SQL. \\
Contenedor & Una «cajita» de Docker que trae un programa listo para usar (PHP, MySQL\ldots) sin instalarlo en tu PC. \\
\texttt{docker-compose.yml} & El «plano» que dice qué contenedores levantar y cómo se conectan. \\
\texttt{echo} & La instrucción que escribe algo en la página. \\
GET & Forma de enviar datos en la URL (\texttt{?nombre=Ana}). Se ven y se pueden guardar en favoritos. \\
\texttt{htmlspecialchars} & El escudo: convierte \texttt{<} y \texttt{>} en texto inofensivo, para que el navegador los muestre en lugar de ejecutarlos. \\
MySQL & El gestor de base de datos: guarda los datos en tablas aunque se apague todo (la \emph{despensa}). \\
PDO & La forma de PHP para conectarse a una base de datos (la \emph{llave} de la despensa). \\
PHP & Lenguaje que se ejecuta en el servidor y genera HTML (el \emph{cocinero}). \\
phpMyAdmin & Página web para ver y editar la base de datos sin escribir SQL. \\
POST & Forma de enviar datos en el cuerpo de la petición. No se ven en la URL (pero no están cifrados). \\
PRG & \emph{Post, Redirect, Get}: después de guardar, redirigir para que F5 no repita el envío. \\
Servidor & La máquina (o el contenedor) donde vive tu sitio y donde corre PHP. \\
Sesión & La memoria que PHP guarda para cada visitante (\texttt{\$\_SESSION}). \\
SQL & El idioma de las bases de datos: \texttt{SELECT} (leer), \texttt{INSERT} (guardar), \texttt{UPDATE}, \texttt{DELETE}. \\
Validar & Revisar que los datos que llegan tengan sentido antes de usarlos. \\
Variable & Una cajita con nombre que guarda un dato. En PHP empieza con \texttt{\$}. \\
XSS & \emph{Cross-Site Scripting}: un ataque en el que alguien mete código HTML o JavaScript en tu página a través de un formulario. \\
\end{longtable}

\subsection*{Funciones que usamos hoy}
\renewcommand{\arraystretch}{1.15}
\begin{tabularx}{\linewidth}{>{\ttfamily}l X}
date('H:i') & hora actual \\
strtoupper(\$t) & a mayúsculas \\
strlen(\$t) & largo del texto \\
trim(\$t) & quita espacios \\
count(\$a) & elementos del array \\
array\_rand(\$a) & posición al azar \\
filter\_var(\$c, \ldots) & valida un correo \\
header('Location: x') & redirige \\
\end{tabularx}

% =====================================================================
\clearpage
\renewcommand{\seccionactual}{APÉNDICE C · RESPUESTAS}
\section*{Apéndice C · Respuestas de «Comprobá que entendiste»}
\addcontentsline{toc}{section}{Apéndice C · Respuestas}

\subsection*{Nivel 0 · El gancho}
\begin{enumerate}
  \item Porque PHP se ejecuta en el servidor: el navegador solo recibe el HTML que PHP generó.
  \item «¡Buenas noches!»: 20 no es menor que 12 ni menor que 19, así que entra al \texttt{else}.
  \item No. El código PHP nunca sale del servidor; solo sale su resultado.
\end{enumerate}

\subsection*{Nivel 1 · Hablar}
\begin{enumerate}
  \item Con el signo \texttt{\$}.
  \item Con simples: \texttt{Total: \$precio} (literal). Con dobles: \texttt{Total: 12.5} (el valor).
  \item Para pegar (concatenar) textos y variables.
  \item Un atajo de \texttt{<?php echo \$x; ?>}.
\end{enumerate}

\subsection*{Nivel 2 · Decidir y repetir}
\begin{enumerate}
  \item \texttt{\$edad < 18}: no es menor que 13, pero sí menor que 18. Sale «Sos adolescente».
  \item \texttt{=} guarda un valor en una variable; \texttt{==} compara dos valores.
  \item Repite un bloque una vez por cada elemento del array. Se escribe una fila y PHP la repite por cada plato.
  \item \texttt{\$item['precio']}.
\end{enumerate}

\subsection*{Nivel 3 · Escuchar}
\begin{enumerate}
  \item Un array con los datos que vinieron en la URL. Lo llena PHP solo.
  \item El \texttt{name} del campo es la clave con la que el dato aparece en \texttt{\$\_GET}.
  \item Porque la contraseña queda en la URL: en el historial, en los registros del servidor y a la vista.
  \item No. POST solo la saca de la URL. Lo que la protege en el camino es HTTPS.
\end{enumerate}

\subsection*{Nivel 4 · Recordar}
\begin{enumerate}
  \item Porque PHP y MySQL están en contenedores distintos. Para el contenedor de PHP, \texttt{localhost} es él mismo; a MySQL lo encuentra por su nombre de servicio, \texttt{db}.
  \item \texttt{query} manda el SQL tal cual (sirve cuando no hay datos del usuario). \texttt{prepare} + \texttt{execute} mandan el SQL con huecos y los datos por separado: es lo seguro cuando hay datos del usuario.
  \item Huecos donde van los datos. MySQL los trata siempre como datos, nunca como parte de la orden SQL.
  \item Para que al recargar (F5) el navegador no vuelva a enviar el formulario y no se guarde el mensaje dos veces.
\end{enumerate}

\subsection*{Nivel 5 · El museo}
\begin{enumerate}
  \item En la línea de arriba (la 5). Muchas veces el error empieza antes, por ejemplo con un \texttt{;} que falta.
  \item Porque no muestra ningún error: la página «funciona», pero muestra algo equivocado. Hay que mirar con atención.
  \item Porque Apache solo le pasa a PHP los archivos que terminan en \texttt{.php}. Los \texttt{.html} los entrega tal cual.
\end{enumerate}

% =====================================================================
\clearpage
\renewcommand{\seccionactual}{APÉNDICE D · SI ALGO FALLA}
\section*{Apéndice D · Si algo falla}
\addcontentsline{toc}{section}{Apéndice D · Si algo falla}

\renewcommand{\arraystretch}{1.4}
\begin{longtable}{>{\raggedright\arraybackslash\bfseries}p{4.6cm} >{\raggedright\arraybackslash}p{11.2cm}}
\toprule
Problema & Qué hacer \\
\midrule
\endhead
El script no se ejecuta («running scripts is disabled») & Primero: \texttt{Set-ExecutionPolicy -Scope Process Bypass}. Vale solo para esa ventana de PowerShell. \\
«Cannot connect to the Docker daemon» & Docker Desktop no está abierto o todavía está arrancando. Abrilo y esperá a que la ballena se quede quieta. \\
«port is already allocated» & Otro programa usa ese puerto (por ejemplo XAMPP en el 80 o un MySQL en el 3306). En \texttt{docker-compose.yml} cambiá el número de la izquierda: \texttt{"8095:80"}. \\
La página dice «La despensa está cerrada» & MySQL todavía está arrancando. Mirá \texttt{docker compose ps}: \texttt{fogon-db} tiene que decir \texttt{healthy}. Esperá unos segundos y recargá. \\
Edité el archivo y no cambia nada & ¿Guardaste (\texttt{Ctrl+S})? ¿Estás editando el archivo de la carpeta \texttt{www} del laboratorio y no una copia? Probá \texttt{Ctrl+F5}. \\
Pantalla en blanco o error de PHP & Leé el mensaje: tipo y línea (ver el nivel 5). Mirá también la línea de arriba. \\
Quiero volver la base de datos a como estaba & \texttt{.\textbackslash scripts\textbackslash reiniciar-base.ps1}. \textbf{Borra} todo lo que guardaste y vuelve a cargar \texttt{init.sql}. \\
Para ver qué dice el servidor & \texttt{docker compose logs web} (PHP y Apache) o \texttt{docker compose logs db} (MySQL). \\
Docker no anda en mi PC & Plan B: \texttt{.\textbackslash scripts\textbackslash plan-b-sin-base.ps1} levanta solo PHP (niveles 0 a 3 y el museo). Si tampoco, trabajá en pareja con un compañero. \\
Apagar todo al final & \texttt{.\textbackslash scripts\textbackslash apagar.ps1}. Los datos de MySQL se conservan para la próxima. \\
\bottomrule
\end{longtable}

\begin{terminal}[PowerShell]{comandos útiles (dentro de la carpeta operacion-fogon)}
docker compose ps              # qué contenedores están corriendo
docker compose logs -f web     # ver los mensajes de PHP en vivo
docker compose restart web     # reiniciar solo PHP
docker compose down            # apagar (los datos quedan)
docker compose down -v         # apagar y BORRAR la base
\end{terminal}

\vfill
\begin{center}
\begin{tikzpicture}
  \fill[negro,rounded corners=3mm] (0,0) rectangle (\linewidth,2.6);
  \fill[fuego,rounded corners=3mm] (0,0) rectangle (0.25,2.6);
  \node[anchor=west,text=oroneon,font=\sffamily\bfseries\large] at (0.8,1.85) {¿Llegaste hasta acá?};
  \node[anchor=west,text=texto,font=\sffamily\small,text width=14.5cm] at (0.8,0.95) {Hoy aprendiste lo mismo que usa el 70\,\% de la web. Mañana lo atacamos\ldots\ y lo blindamos. Traé tu proyecto.};
\end{tikzpicture}
\end{center}
